\documentclass[11pt,a4paper]{article}
\usepackage[margin=25mm]{geometry}
\usepackage{amsmath,amssymb,amsthm,hyperref}
\setlength{\emergencystretch}{2em}
\newtheorem{lemma}{Lemma}
\newtheorem{proposition}[lemma]{Proposition}
\newcommand{\gen}[1]{\langle#1\rangle}
\title{An internal block action in the remaining support branch}
\author{Human Compression project: bounded L044 proof component}
\date{}
\begin{document}\maketitle
Let $G=\gen{\rho_0,\ldots,\rho_{r-1}}\leq A_n$ be transitive,
$n\geq12$, with a connected string C-marking and every maximal
omission intransitive. Let the unique eligible $i$-edge
$\{a,b\}$, $0<i<r-1$, join the two $G_i$-orbits $A,B$.
Assume $|B|\geq2$ and $L=G_{<i}$ fixes $B$ pointwise and is intransitive on $A$.
Write $\beta=a^L$, $k=|\beta|$, and $\tau=\rho_{i+1}$.

The $L$-orbits in $A$ are blocks for $G_i=L G_{>i}$, because
the two factors commute. They have equal sizes. Every
$\rho_j$ with $j\geq i+2$ fixes $a$ and centralizes $L$, so
fixes $\beta$ pointwise. Thus $\tau$ must send $\beta$ to
a different block $\beta'$, otherwise $\beta$ would be
$G_i$-invariant. Every generator of index at least $i+3$
also fixes $\beta'$ pointwise, by commutation with $\tau$.
The involution $\tau$ moves $b$ as well: it is the only
member of $G_i$ that can do so, and $B$ is a nontrivial $G_i$-orbit.

\begin{lemma}[A faithful commuting-orbit action]
The action of $L$ on each of its orbits in $A$ is faithful,
and the actions on $\beta$ and $\beta'$ are identified by
$\tau$. Suppose $\sigma=\rho_{i+2}$ exists, preserves
$\beta'$, and acts nontrivially there. Then $A=\beta\cup\beta'$,
$k$ is even, and
\[
                              i\leq k/2-1.               \tag{1}
\]
\end{lemma}
\begin{proof}
A member of $L$ fixing one $L$-orbit pointwise fixes all its
translates under $G_{>i}$, by commutation. Those translates
cover $A$; it also fixes $B$, hence is the identity. This
proves faithfulness and the asserted identification.

The set $\beta\cup\beta'$ is preserved by $L$, interchanged
blockwise by $\tau$, preserved by $\sigma$, and fixed
pointwise by every later generator. It is therefore all of
$A$. The restriction of $\sigma$ to $\beta'$ centralizes
the transitive $L$-action. Being nontrivial and involutory,
it has no fixed point, so supplies $k/2$ pairs. The conjugate
$\sigma^\tau$ supplies corresponding pairs on $\beta$.
These are $L$-invariant block systems.

Choose a minimal distinguished subset $M$ of $L$ generating
its whole action on these $k/2$ pairs. This block action is
transitive, so $\gen M$ has at most two orbits in $\beta$.
If there are two, $\sigma^\tau$ interchanges them; hence
$\gen{M,\sigma^\tau}$ is transitive on $\beta$ in either
case. Similarly $\gen{M,\sigma}$ is transitive on $\beta'$.
Together with $\tau$ these groups are transitive on $A$.

If a distinguished $L$ generator $\rho_j$ were outside $M$,
all the displayed elements would belong to $G_j$, so $G_j$
would be transitive on $A$. But $\rho_j$ is supported on $A$,
and could not join distinct $G_j$-orbits. This contradicts
transitivity of $G=\gen{G_j,\rho_j}$ and intransitivity of
$G_j$. Thus every $L$ generator belongs to $M$. Their pair
images are an independent family in $S_{k/2}$; the retained
Whiston bound gives (1). Since $i>0$, necessarily $k\geq4$.
No projected intersection property is used.
\end{proof}

\begin{proposition}[Strict support deficit or completed rank bound]
Under the internal-block hypothesis of the lemma, if
$i+2=r-1$, then $r\leq\lfloor(n-1)/2\rfloor$.
Otherwise $h=i+2$ is the largest shared label, and with
$X=\operatorname{supp}(G_{>h})$ one has
\[
 X\subseteq B\setminus\{b\},\qquad
 |\operatorname{Fix}(G_{>h})|\geq2k+1,\qquad
 h\leq\frac{n-|X|-3}{2}.
\]
In particular the equality support configuration in source
Proposition~5.23 cannot arise in this subcase.
\end{proposition}
\begin{proof}
If $i+2=r-1$, (1) and $n\geq2k+2$ give
\[
 r=i+3\leq k/2+2\leq k\leq\left\lfloor\frac{n-1}{2}\right\rfloor,
\]
using $k\geq4$.

If a later generator exists, $\rho_{i+3}$ fixes $A$ and acts
nontrivially on $B$. Were $\sigma$ trivial on $B$, their
supports would be disjoint, so these adjacent generators
would commute. All generators on opposite sides of this
adjacent pair already commute by the string condition;
this would disconnect the noncommutation diagram. Hence
$\sigma$ also acts nontrivially on $B$. It is the last
shared generator because all later generators fix $A$.
The tail also fixes $b$, by the unique-crossing lemma.
It fixes at least $2k+1$ points, while (1) gives
\[
                        2h+3=2i+7\leq k+5\leq2k+1
\]
for $k\geq4$. This proves the strict bound and the support
placement claimed.
\end{proof}

\paragraph{Scope and attribution.}
This is the internal-block subcase of Cameron--Fernandes--Leemans--Mixer,
\emph{Highest rank of a polytope for $A_n$}, arXiv:1605.09173v1,
Proposition~5.23, physical page~32. The faithful commuting-orbit
and actual-omission argument makes its rank reduction explicit,
using the paid unique-crossing proof and retained Whiston bound.
The application to $G=A_n$ has a connected diagram by simplicity;
for a general subgroup the connectedness assumption is retained.
The nontrivial-$B$ hypothesis is explicit in this successor: being the
only possible mover of $b$ alone does not exclude a singleton orbit.
A shared generator in the intended source branch ensures $|B|\geq2$.
No finite computation or formal theorem is added. The case where
$\rho_{i+2}$ moves $\beta'$ to another block, all further
path/equality rigidity and the whole L044 contract remain separate.
\end{document}
